\documentclass[11pt,a4paper]{article}
\usepackage[utf8]{inputenc}
\usepackage[T1]{fontenc}
\usepackage[portuguese]{babel}
\usepackage{geometry}
\usepackage{helvet}
\usepackage{enumitem}
\usepackage{titlesec}
\usepackage[table]{xcolor}
\usepackage{hyperref}
\usepackage{tabularx}
\usepackage{booktabs}
\usepackage{fancyhdr}
\usepackage[most]{tcolorbox}
\usepackage{multicol}
\usepackage{ragged2e}

\geometry{margin=2.2cm, top=2.5cm, bottom=2.5cm}
\renewcommand{\familydefault}{\sfdefault}

% ── Paleta de cores ──
\definecolor{uefsdark}{RGB}{0,51,102}
\definecolor{uefsblue}{RGB}{25,84,150}
\definecolor{uefslight}{RGB}{232,241,250}
\definecolor{uefsgray}{RGB}{110,110,110}
\definecolor{uefsaccent}{RGB}{180,60,40}

% ── Seções ──
\titleformat{\section}
  {\normalfont\large\bfseries\color{uefsdark}\scshape}{}{0em}{}%
  [\vspace{2pt}{\color{uefsdark}\titlerule}]
\titlespacing{\section}{0pt}{20pt}{8pt}

\titleformat{\subsection}
  {\normalfont\normalsize\bfseries\color{uefsblue}}{}{0em}{}
\titlespacing{\subsection}{0pt}{12pt}{4pt}

% ── Listas ──
\setlist[itemize]{leftmargin=1.5em,topsep=4pt,itemsep=2pt,parsep=0pt}
\setlist[enumerate]{leftmargin=2em,topsep=4pt,itemsep=4pt,parsep=0pt}

% ── tcolorbox ──
\tcbset{
  uefsbox/.style={
    colback=uefslight, colframe=uefsdark,
    boxrule=0.4pt, arc=3pt,
    left=8pt, right=8pt, top=6pt, bottom=6pt,
    fontupper=\small
  },
  ementabox/.style={
    colback=white, colframe=uefsdark,
    boxrule=0.8pt, arc=0pt,
    left=10pt, right=10pt, top=8pt, bottom=8pt,
    title={\scshape\bfseries Ementa},
    fonttitle=\large\color{white},
    coltitle=white, colbacktitle=uefsdark,
    attach boxed title to top left={yshift=-2mm,xshift=2mm},
    boxed title style={arc=2pt,boxrule=0pt}
  }
}

% ── Cabeçalho / Rodapé ──
\pagestyle{fancy}
\fancyhf{}
\renewcommand{\headrulewidth}{0.4pt}
\renewcommand{\headrule}{\hbox to\headwidth{\color{uefsdark}\leaders\hrule height \headrulewidth\hfill}}
\fancyhead[L]{\small\color{uefsgray}UEFS\,--\,PLANTERR}
\fancyhead[R]{\small\color{uefsgray}PLA001\,--\,Plano de Ensino 2026.1}
\fancyfoot[C]{\small\color{uefsgray}\thepage}

\hypersetup{colorlinks=true,linkcolor=uefsblue,urlcolor=uefsblue}

% ══════════════════════════════════════════════════════════
\begin{document}
\thispagestyle{empty}

% ── Capa / Cabeçalho institucional ──
\begin{center}
\color{uefsdark}\rule{\textwidth}{1.2pt}\\[0.35cm]
{\Large\bfseries UNIVERSIDADE ESTADUAL DE FEIRA DE SANTANA}\\[0.20cm]
{\normalsize Departamento de Ciências Humanas e Filosofia}\\[0.08cm]
{\normalsize Programa de Pós-Graduação em Planejamento Territorial\,--\,PLANTERR}\\[0.20cm]
\color{uefsdark}\rule{\textwidth}{1.2pt}
\end{center}

\vspace{0.8cm}

\begin{center}
{\LARGE\bfseries\color{uefsdark} PLA001\,--\,Métodos e Técnicas para Elaboração de\\Projetos de Pesquisa e Intervenção}\\[0.35cm]
{\large (Teórica-Prática\,--\,01)}\\[0.60cm]
{\Large\bfseries Plano de Ensino\,--\,2026.1}\\[0.30cm]
{\normalsize 15 encontros\,$\times$\,4\,h\enspace|\enspace Quintas-feiras\enspace|\enspace 07h30\,--\,11h30}\\[0.80cm]
{\normalsize\textbf{Professores:} Prof.\,Dr.\,Luiz Diego Vidal \enspace\&\enspace Prof.\,Dr.\,Danilo Uzeda da Cruz}
\end{center}

\vspace{0.6cm}

% ══════════════════════════════════════════════════════════
\section*{Objetivo Geral}

Capacitar os discentes do Mestrado Profissional em Planejamento Territorial na elaboração, condução e avaliação de projetos de pesquisa e intervenção territorial, integrando fundamentos epistemológicos, metodologias qualitativas e quantitativas, e práticas participativas para o diagnóstico e a transformação de realidades socioambientais complexas.

% ══════════════════════════════════════════════════════════
\section*{Objetivos Específicos}

\begin{enumerate}
\item Compreender os fundamentos epistemológicos da pesquisa científica e sua aplicação ao planejamento territorial, reconhecendo diferentes paradigmas metodológicos e suas implicações éticas e políticas.
\item Distinguir projetos de pesquisa acadêmica e projetos de intervenção social, identificando suas especificidades metodológicas, desafios operacionais e contribuições para a transformação territorial.
\item Formular problemas de pesquisa e questões de intervenção territorialmente contextualizados, alinhados a demandas sociais e ambientais concretas.
\item Aplicar metodologias qualitativas (entrevistas, diagnóstico participativo, caminhada transversal, cartografia social) na construção de diagnósticos territoriais participativos e culturalmente sensíveis.
\item Utilizar métodos quantitativos (estatística descritiva e inferencial, testes de hipóteses, análise de correlação e regressão) para análise de dados ambientais e territoriais.
\item Integrar abordagens qualitativas e quantitativas em desenhos de pesquisa mistos, potencializando a compreensão de fenômenos territoriais complexos.
\item Planejar, executar e sistematizar atividades de campo em contextos reais de intervenção territorial, articulando teoria e prática.
\item Avaliar criticamente projetos de políticas públicas e intervenções territoriais quanto à sua fundamentação metodológica, coerência interna e impactos socioambientais.
\item Desenvolver postura ética, reflexiva e crítica diante dos dilemas da pesquisa aplicada e da intervenção social em contextos de desigualdade e conflito socioambiental.
\item Comunicar resultados de pesquisa de forma clara, rigorosa e acessível a diferentes públicos (acadêmico, técnico, comunitário, gestor público).
\end{enumerate}

% ══════════════════════════════════════════════════════════
\section*{Programa do Componente Curricular}

\subsection*{Unidade\,I\,--\,Fundamentos da Pesquisa Aplicada ao Planejamento Territorial\enspace{\normalfont\small(Enc.\,01--08)}}
\begin{itemize}
\item Pesquisa científica: princípios epistemológicos, etapas e importância
\item Projetos de intervenção e ciências sociais: metodologias e dilemas éticos
\item Produção do conhecimento e transformação social
\item Método científico e intervenção territorial
\item Análise de projetos em políticas sociais, gestão governamental e política ambiental
\end{itemize}

\subsection*{Unidade\,II\,--\,Metodologias de Pesquisa e Análise de Dados Territoriais\enspace{\normalfont\small(Enc.\,09--15)}}
\begin{itemize}
\item Metodologias qualitativas: entrevistas, diagnóstico participativo, caminhada transversal e cartografia social
\item Metodologias quantitativas: estatística descritiva, medidas de tendência central, inferência estatística, correlação, regressão e testes de hipóteses aplicados a dados ambientais
\item Pesquisa de campo: planejamento, execução e sistematização de dados primários
\item Integração quali-quantitativa na análise territorial
\end{itemize}

% ══════════════════════════════════════════════════════════
\section*{Significado do Componente para a Formação Profissional}

A disciplina constitui um componente estratégico para a formação de profissionais no Mestrado Profissional em Planejamento Territorial, atendendo às demandas contemporâneas por abordagens integradas, baseadas em evidências e comprometidas com a sustentabilidade socioambiental.

\subsection*{Relevância para o Planejamento Territorial Contemporâneo}
O planejamento territorial enfrenta desafios cada vez mais complexos, que demandam capacidade analítica para integrar múltiplas dimensões (ambiental, social, econômica, institucional) e escalas (local, regional, nacional). Esta disciplina fornece instrumentos metodológicos essenciais para que profissionais possam: (i)~diagnosticar problemas territoriais de forma rigorosa e sistemática, (ii)~fundamentar decisões em dados empíricos confiáveis, (iii)~avaliar impactos de políticas e projetos, e (iv)~comunicar resultados de forma transparente e acessível.

\subsection*{Articulação entre Teoria e Prática}
Ao combinar fundamentos epistemológicos com metodologias aplicadas e práticas de campo, a disciplina prepara profissionais para atuar na interface entre academia, gestão pública e organizações da sociedade civil. A ênfase em metodologias participativas reflete a crescente valorização do protagonismo comunitário e dos saberes locais nos processos de planejamento, alinhando-se aos princípios da gestão democrática do território.

\subsection*{Competências para Atuação Profissional}
O domínio de métodos qualitativos e quantitativos de análise de dados ambientais é cada vez mais exigido em concursos públicos, processos seletivos de organizações não governamentais e consultoria técnica especializada. Profissionais capazes de conduzir diagnósticos territoriais, avaliar políticas públicas e propor intervenções baseadas em evidências possuem diferencial competitivo significativo no mercado de trabalho.

\subsection*{Contribuição para a Dimensão Ambiental do Planejamento}
Em contexto de crise climática, perda de biodiversidade e agudização de conflitos socioambientais, a capacidade de analisar dados ambientais de forma crítica e contextualizada é indispensável. Esta disciplina capacita profissionais a interpretar indicadores de qualidade ambiental, avaliar impactos de intervenções territoriais e propor soluções que conciliem desenvolvimento socioeconômico e conservação ambiental.

\subsection*{Ética e Responsabilidade Social}
A disciplina promove reflexão crítica sobre dilemas éticos e políticos da pesquisa aplicada, formando profissionais conscientes de sua responsabilidade social e comprometidos com processos de planejamento territorial justos, inclusivos e ambientalmente sustentáveis.

% ══════════════════════════════════════════════════════════
\section*{Cronograma}

\begin{small}
\renewcommand{\arraystretch}{1.35}
\begin{tabularx}{\textwidth}{@{}c c X@{}}
\toprule
\textbf{Enc.} & \textbf{Data} & \textbf{Tema} \\
\midrule
01 & 12/03\,(Qui) & Apresentação do curso \\
02 & 19/03\,(Qui) & Pesquisa científica: princípios, etapas e importância\,(I) \\
03 & 26/03\,(Qui) & Pesquisa científica: princípios, etapas e importância\,(II) \\
04 & 02/04\,(Qui) & Projetos de intervenção e as ciências sociais: metodologias e dilemas \\
05 & 09/04\,(Qui) & O problema de pesquisa em um projeto de intervenção social -- Produção do conhecimento humano $\times$ conhecimento como fator de (re)produção na sociedade \\
06 & 16/04\,(Qui) & O método científico e a intervenção social: o que é um projeto de intervenção social? \\
\rowcolor{uefslight}
   & \textit{21/04\,(Ter)} & \textit{Feriado -- Tiradentes (não afeta quintas)} \\
07 & 23/04\,(Qui) & Seminário\,A -- Projetos: Bloco\,1 (Políticas Sociais) \\
08 & 30/04\,(Qui) & Seminário\,B -- Projetos: Bloco\,2 (Gestão Governamental) \\
09 & 07/05\,(Qui) & Metodologias qualitativas: entrevista estruturada, semiestruturada e aberta; Diagnóstico Participativo e caminhada transversal; Cartografia social (Proj.\,Quilombo Legal) \\
10 & 14/05\,(Qui) & Seminário\,C -- Projetos: Bloco\,3 (Política Ambiental) \\
11 & 21/05\,(Qui) & Metodologias Quantitativas\,I: introdução à estatística, tipos de variáveis, estatística descritiva e medidas de tendência central \\
12 & 28/05\,(Qui) & Metodologias Quantitativas\,II: inferência estatística, correlação, regressão e testes de hipóteses + Preparação para atividade de campo \\
\rowcolor{uefslight}
13 & 02--03/06\,(Seg--Ter) & \textbf{Atividade de campo} -- Projetos sociais e políticas públicas (dois dias, fora do calendário regular) \\
14 & 04/06\,(Qui) & Sistematização dos dados de campo e análise preliminar \\
15 & 11/06\,(Qui) & Apresentação dos resultados + Avaliação e Encerramento do Semestre \\
\bottomrule
\end{tabularx}
\end{small}

\begin{tcolorbox}[uefsbox]
\textbf{Estrutura geral:}\enspace
Enc.\,01--08 → Bloco Teórico-Metodológico\enspace|\enspace
Enc.\,09--15 → Bloco Qualitativo-Quantitativo\,+\,Campo\,+\,Avaliação\enspace|\enspace
Enc.\,13: dois dias consecutivos (Seg--Ter).
\end{tcolorbox}

% ══════════════════════════════════════════════════════════
\section*{Metodologia Geral}

A disciplina adota uma abordagem híbrida que integra fundamentos teóricos, metodologias aplicadas e práticas de campo, estruturada em três eixos pedagógicos complementares:

\subsection*{1.\enspace Aulas Expositivas Dialogadas}
Apresentação dos fundamentos epistemológicos da pesquisa científica, dos métodos de intervenção social e das técnicas de análise de dados ambientais, promovendo debates sobre dilemas éticos e metodológicos no planejamento territorial.

\subsection*{2.\enspace Seminários Temáticos}
Os discentes organizarão e apresentarão projetos em três blocos (Políticas Sociais, Gestão Governamental e Política Ambiental), exercitando a análise crítica de intervenções territoriais reais, com ênfase em metodologias participativas e na articulação entre dimensões sociais, ambientais e institucionais.

\subsection*{3.\enspace Práticas de Campo e Análise de Dados}
Aplicação direta das metodologias qualitativas (entrevistas, diagnóstico participativo, cartografia social) e quantitativas (estatística descritiva e inferencial) em contexto real, seguida de sistematização coletiva dos dados coletados, promovendo o aprendizado experiencial e a reflexão sobre desafios práticos da pesquisa aplicada ao território.

\vspace{6pt}
A avaliação será processual e contínua, baseada em: (i)~participação ativa nos debates, (ii)~qualidade dos seminários apresentados, (iii)~relatório da atividade de campo, e (iv)~capacidade de articulação teórico-prática demonstrada ao longo do semestre.

% ── Avaliação ──
\begin{tcolorbox}[uefsbox, colframe=uefsaccent, colback=white, title={\bfseries\color{white}\scshape Avaliação}, fonttitle=\normalsize, coltitle=white, colbacktitle=uefsaccent, attach boxed title to top left={yshift=-2mm,xshift=2mm}, boxed title style={arc=2pt,boxrule=0pt}]
\begin{itemize}[leftmargin=1.2em]
\item Participação em debates e seminários\,\dotfill\,\textbf{30\,\%}
\item Apresentação de seminários temáticos\,\dotfill\,\textbf{30\,\%}
\item Relatório de atividade de campo\,\dotfill\,\textbf{40\,\%}
\end{itemize}
\end{tcolorbox}

% ══════════════════════════════════════════════════════════
\section*{Habilidades e Competências}

\begin{multicols}{2}
\subsection*{Competências Metodológicas}
\begin{itemize}
\item Dominar os fundamentos epistemológicos da pesquisa científica aplicada ao planejamento territorial
\item Selecionar e aplicar metodologias qualitativas e quantitativas adequadas aos diferentes contextos de intervenção territorial
\item Integrar múltiplas fontes de dados (ambientais, sociais, institucionais) na análise de problemas territoriais complexos
\item Utilizar ferramentas estatísticas para análise descritiva e inferencial de dados ambientais
\end{itemize}

\subsection*{Competências Analíticas e Críticas}
\begin{itemize}
\item Avaliar criticamente projetos de intervenção social e políticas públicas territoriais
\item Identificar e analisar dilemas éticos e metodológicos em contextos de planejamento participativo
\item Articular conhecimento teórico-conceitual com práticas de diagnóstico e intervenção territorial
\item Interpretar indicadores ambientais e sociais para subsidiar processos decisórios
\end{itemize}

\columnbreak

\subsection*{Competências Práticas e Relacionais}
\begin{itemize}
\item Conduzir entrevistas estruturadas, semiestruturadas e abertas com diferentes atores sociais
\item Aplicar metodologias participativas (diagnóstico participativo, caminhada transversal, cartografia social)
\item Trabalhar em equipe multidisciplinar na construção de diagnósticos territoriais
\item Comunicar resultados de pesquisa de forma clara e objetiva para públicos técnicos e comunitários
\end{itemize}

\subsection*{Competências Estratégicas}
\begin{itemize}
\item Formular problemas de pesquisa alinhados a demandas de intervenção territorial
\item Desenhar projetos de intervenção social fundamentados em evidências empíricas
\item Integrar conhecimento científico e saberes locais na construção de soluções territoriais
\item Posicionar-se como profissional reflexivo e agente de transformação territorial
\end{itemize}
\end{multicols}

% ══════════════════════════════════════════════════════════
\begin{tcolorbox}[ementabox]
Pesquisa científica: princípios, etapas e importância. A relação entre pesquisa e intervenção na produção científica. Projetos de intervenção: a ação planejada e teoricamente embasada. Produção do conhecimento científico e tecnológico. Modelos de Pesquisa Social, Tecnológica e de Planejamento Territorial.
\end{tcolorbox}

% ══════════════════════════════════════════════════════════
\section*{Bibliografia Básica}

\begin{small}
\hangindent=1.5em\hangafter=1
ARMANI, D. \textbf{Como elaborar projetos?}: guia prático para elaboração e gestão de projetos sociais. Porto Alegre: Tomo Editorial, 2009.\par\medskip

\hangindent=1.5em\hangafter=1
BECKER, H.\,S. \textbf{De que lado estamos?} Uma teoria da ação coletiva. Rio de Janeiro: Zahar, 1977.\par\medskip

\hangindent=1.5em\hangafter=1
BRACAGIOLI NETO, A.; GEHLEN, I.; OLIVEIRA, V.\,L. \textbf{Planejamento e gestão de projetos para o desenvolvimento rural}. Porto Alegre: UFRGS, 2010.\par\medskip

\hangindent=1.5em\hangafter=1
BROSE, M. (Org.). \textbf{Metodologia participativa}: uma introdução a 29 instrumentos. Porto Alegre: Tomo Editorial, 2001.\par\medskip

\hangindent=1.5em\hangafter=1
BUARQUE, S.\,C. \textbf{Metodologia de planejamento do desenvolvimento sustentável}. Brasília: IICA, 1995.\par\medskip

\hangindent=1.5em\hangafter=1
CRUZ, D.\,U. \textbf{Planejamento participativo e políticas públicas}: participação social e metodologias participativas no Brasil contemporâneo. Feira de Santana: Z Arte, 2016.\par\medskip

\hangindent=1.5em\hangafter=1
HAGUETTE, T.\,M.\,F. \textbf{Metodologias Qualitativas em Sociologia}. 4.\,ed. Rio de Janeiro: Vozes, 2013. p.\,109--170.\par\medskip

\hangindent=1.5em\hangafter=1
PEREIRA, M.\,N. \textit{et al.} \textbf{Métodos e Meios de Comunicação em Extensão Rural}. Porto Alegre, 2009. Disponível em: \url{http://www.emater.tche.br/site/arquivos_pdf/teses/METODOSDEEXTENSAOGLOSSARIO.pdf}. Acesso em: 20 jun. 2022.\par\medskip

\hangindent=1.5em\hangafter=1
THIOLLENT, M. \textbf{Metodologia da pesquisa-ação}. 14.\,ed. São Paulo: Cortez, 2005.\par\medskip

\hangindent=1.5em\hangafter=1
UZÊDA, L.\,F.\,F.; CRUZ, D.\,U. \textbf{Extensão rural no Brasil}: percursos, metodologias e desafios. Salvador: Pinaúna, 2020.\par\medskip

\hangindent=1.5em\hangafter=1
YIN, R.\,K. \textbf{Estudo de Caso}: planejamento e métodos. Porto Alegre: Bookman, 2001.
\end{small}

\end{document}
